\begin{proof}
\noindent\textbf{Step 1.} Choose coordinates $(x,y,z)$ with critical curve
$\{x=y=0\}$. The rank-one quadratic form along this curve
can be written
\[
 a(z)x^2.
\]
To obtain these coordinates without a meromorphic change,
divide the matrix of its coefficients by its common
vanishing order at $z=0$. A primitive symmetric rank-one
matrix has a unit entry, factors as a unit times the
square of a primitive linear form, and that linear form
extends to an invertible holomorphic change of the normal
coordinates. The intrinsic cubic coefficient on its
kernel is then $b(z)y^3$. The generic cusp hypothesis says
that both $a$ and $b$ are nonzero functions; shrink the
disc so that their only possible zero is at $z=0$.

\noindent\textbf{Step 2.} If $a(0)\ne0$, the Morse lemma with parameters gives
\[
 f=X^2+y^3u(y,z),\qquad u(0,z)=b(z)
\]
up to multiplication of $b$ by a holomorphic unit.
If $u(0,0)=0$, the reduced support of $y^3u(y,z)$ is
singular: $u$ has a component through the origin
different from $y=0$, because $u(0,z)$ is not identically
zero. The arbitrary-suspension clause of \ref{cdp:st:A.5.17} then
puts a sphere of nonzero self-intersection into a
nearby actual torus fibre, which is impossible.
Thus $u$ is a unit, and replacing $y$ by a holomorphic
cube root of $u$ times $y$ gives \eqref{cdp:eq:C.CR.15}.

\noindent\textbf{Step 3.} It remains to exclude $a(0)=0$. Consider the actual
holomorphic perturbation
\[
                  f_{\epsilon,\delta}
                      =f+\epsilon x^2+\delta y^3.
 \tag{C.CR.16}\label{cdp:eq:C.CR.16}
\]
Choose small parameters so that $a(z)+\epsilon$ has a
simple zero $z_\epsilon$ and
$b(z_\epsilon)+\delta\ne0$. Such choices exist:
regular values of the nonconstant function $a$ are
dense near zero, and finitely many values of $\delta$
are then excluded.

\noindent\textbf{Step 4.} We specify the germ produced at
$(0,0,z_\epsilon)$. Weierstrass preparation in $y$,
followed by a Tschirnhaus substitution, writes it as
a holomorphic unit times
\[
             Y^3+x^2\bigl(\beta(x,z)Y+\gamma(x,z)\bigr).
 \tag{C.CR.17}\label{cdp:eq:C.CR.17}
\]
Here the absence of a linear transverse jet gives
the divisibility of the constant term by $x^2$,
and the rank-one quadratic form gives the divisibility
of the coefficient of $Y$ by $x^2$. More explicitly,
before the Tschirnhaus substitution the coefficient
of $y^2$ is divisible by $x$, while those of $y$
and of $1$ are divisible by $x^2$; these divisibilities
are preserved by that substitution.
The function $\gamma(0,z)$ is a unit multiple of
$a(z)+\epsilon$. Consequently
\[
             Z=\beta(x,z)Y+\gamma(x,z)
\]
is an invertible change of the curve coordinate at
the chosen simple zero. Equation \eqref{cdp:eq:C.CR.17} is thus
a unit times $Y^3+x^2Z$. This remaining unit can
also be removed: this polynomial is homogeneous
of degree three, so scaling all three coordinates
by a holomorphic cube root of the unit gives the
exact equation
\[
                            Y^3+x^2Z.
 \tag{C.CR.18}\label{cdp:eq:C.CR.18}
\]
\ref{cdp:st:A.5.15} supplies a compact spherical cycle of square $-4$
in an arbitrarily small nonzero level of this germ.
Indeed its Hirzebruch formula is
$-2(b/\gcd(b,c))(c/\gcd(b,c))$, here with
$(b,c)=(2,1)$. The assertion is about the homology
class of the explicit continuous map from $S^2$;
no embedded representative is needed.

\noindent\textbf{Step 5.} For completeness we prove the boundary comparison
needed to put this spherical cycle in the original fibre.
Choose a sufficiently small closed ball whose
boundary meets the smooth critical curve transversely
and is transverse to the smooth stratum of $f=0$.
On the compact link of the critical curve, both
$a$ and $b$ are units. For all sufficiently small
$(\epsilon,\delta)$ the parametrized Morse and
cusp-coordinate changes therefore give holomorphic
product-cusp charts there, retaining $z$ as curve
coordinate. In each such chart the vector
$\partial/\partial z$ is tangent to every level
of $f_{\epsilon,\delta}$. Along the critical link,
an appropriate real radial vector in this direction
is transverse to the fixed sphere. This remains
true on a neighbourhood of that link, uniformly
for the small perturbations. Finitely many charts
suffice. On the rest of the boundary ordinary
compactness and transversality at parameter zero
give the same conclusion for small target values.
Thus all sufficiently small nonzero levels have
uniformly transverse boundary. The full-critical-
locus hypothesis ensures that no other boundary
singular curve has been omitted from this argument.

\noindent\textbf{Step 6.} For fixed small parameters, the nonzero critical
values inside the ball are finite. The local
spherical cycle from \eqref{cdp:eq:C.CR.18} can therefore be transported
through regular target values to a fixed small
value $t_1$, avoiding those critical values.
Choose the parameters small enough that every
critical value throughout the segment joining
them to zero has modulus less than $|t_1|/2$.
This is possible by compactness: every limit of
critical points is a critical point of $f$, and
its value is zero; the just-proved product charts
control the boundary as well. The family of
$t_1$-levels along this parameter segment is a
proper submersion with transverse boundary.
Ehresmann transport identifies it with the
original $f=t_1$ level and preserves the square
of the transported spherical class.

\noindent\textbf{Step 7.} The original level lies in a smooth nearby torus
fibre. A two-torus is aspherical, so every sphere
in it is null-homotopic and has zero homology
class and self-intersection. This contradicts
the nonzero square obtained above. Hence $a(0)\ne0$,
and the first part of the proof gives \eqref{cdp:eq:C.CR.15}.
If the original map is a unit times $f^m$, absorb
the unit by its local $m$th root before this
comparison; the resulting sphere lies in the
actual level with parameter $t_1^m$.
\end{proof}







